% year: תשפ"ג
% semester: א
% moed: ג
% source: exams/מבחנים/תשפג/מועד ג תשפג.pdf
% number: 6א
% categories: improper-integrals

\begin{question}
בדקו התכנסות של האינטגרל $\displaystyle\int_0^\infty e^{-x^2}dx$
\end{question}

\begin{hint}
את האינטגרל לא ניתן לחשב בפונקציות אלמנטריות. פצלו ב-$x=1$ והשוו את $e^{-x^2}$ לפונקציה פשוטה יותר עבור $x\ge1$.
\end{hint}

\begin{solution}
נפצל: $\int_0^\infty e^{-x^2}dx=\int_0^1 e^{-x^2}dx+\int_1^\infty e^{-x^2}dx$.
האינטגרל הראשון הוא אינטגרל מסוים רגיל של פונקציה רציפה על קטע סגור, ולכן קיים וסופי.

עבור $x\ge 1$ מתקיים $x^2\ge x$, ולכן $0<e^{-x^2}\le e^{-x}$. האינטגרל
\[ \int_1^\infty e^{-x}dx=\lim_{R\to\infty}\left(e^{-1}-e^{-R}\right)=e^{-1} \]
מתכנס. לפי \textbf{מבחן ההשוואה} לאינטגרלים לא-אמיתיים של פונקציות אי-שליליות, גם $\int_1^\infty e^{-x^2}dx$ מתכנס. לכן $\boxed{\int_0^\infty e^{-x^2}dx\ \text{מתכנס}}$
(וערכו, כידוע, $\frac{\sqrt\pi}{2}$, אך זה לא נדרש).
\end{solution}
